% ============================================================
% Section 2.4 + Proposition 1 (β₁ stripped)
% Drop-in replacement content for main.tex.
% ============================================================

\subsection{Model space and revision graph}\label{sec:model_space}

The preceding sections characterize DLN stages by the within-episode
belief-dependency graph $G$ that an agent uses to represent outcomes, exposures,
and option relations.  Our Network instantiation also implements an explicit
structural learning cycle that can \emph{revise which belief graph is active}
when a structural hypothesis fails.  To formalize this two-level distinction, we
separate (i) inference and learning \emph{within} a belief graph (Level~1) from
(ii) monitoring and revising the belief graph itself (Level~2).

\paragraph{Model space.}
Let $\mathcal{M}$ be a set of candidate belief-dependency structures.  In the
present instantiation, the core model space contains two elements:
a compressed factor graph $G_F$ (shared latent factors connecting options) and
an expanded tabular graph $G_{\mathrm{tab}}$ (independent option beliefs).
Formally, $\mathcal{M}=\{G_F,\,G_{\mathrm{tab}}\}$.

\paragraph{Revision graph.}
Define the \emph{revision graph} as a directed graph over model space,
$\mathcal{R}=(\mathcal{M},\mathcal{T})$, where $\mathcal{T}\subseteq
\mathcal{M}\times\mathcal{M}$ is the set of revision transitions available to an
agent.  In our simulation, the key transitions are:
(i)~\emph{expansion} (switching from $G_F$ to $G_{\mathrm{tab}}$ when the factor
hypothesis fails) and (ii)~\emph{contraction} (switching from $G_{\mathrm{tab}}$
back to $G_F$ when predictive evidence supports compression again).

\paragraph{Completing the cycle: contraction.}
Expansion alone yields an irreversible path $G_F \rightarrow
G_{\mathrm{tab}}$.  Adding a contraction transition
$(G_{\mathrm{tab}},G_F)\in\mathcal{T}$ completes a revision cycle.  This return
transition matters even when $|\mathcal{M}|=2$: it is the mechanism that
enables an agent to recover compressed $O(F)$ operation after a period of
operating in an expanded $O(K)$ representation.

\paragraph{Revision capacity by DLN stage.}
The revision graph characterizes each stage's meta-cognitive capacity in
graph-combinatorial terms:
%
\begin{itemize}[leftmargin=*]
\item \textbf{Dot.}  No model space ($|\mathcal{M}|=0$); $\mathcal{R}$
  is empty.  No revision capacity.
\item \textbf{Linear.}  $\mathcal{M} = \{G_{\mathrm{tab}}\}$; a single
  vertex with no edges.  The agent occupies a fixed point in model space.
\item \textbf{Network (expand-only).}  $\mathcal{M} =
  \{G_F, G_{\mathrm{tab}}\}$; one directed edge
  $G_F \to G_{\mathrm{tab}}$.  The agent can detect model failure and
  expand, but cannot return to a compressed representation.
\item \textbf{Network (full cycle).}  Two directed edges (expand and
  contract).  $\mathcal{R}$ contains a cycle: the agent can depart from a
  model, operate under an alternative, and return.  This is the
  graph-theoretic signature of full revision capacity.
\end{itemize}
%
The functional consequence of the return transition is stated in
Proposition~\ref{prop:crossover}(iii): an agent that can return to
$G_F$ recovers $O(F)$ cost scaling in bounded time; an agent that
cannot is permanently locked at $O(K)$.

\paragraph{Two levels of computation.}
Level~1 computation concerns inference and parameter learning within a chosen
belief graph $G\in\mathcal{M}$ (e.g., updating factor means or option means).
Level~2 computation concerns monitoring model adequacy and executing transitions
in $\mathcal{R}$.  Proposition~\ref{prop:crossover}(iii) establishes that an
agent whose revision graph includes a return transition recovers $\Theta(F)$
scaling in bounded time once factor structure is present, while an agent without
a return transition remains at $\Theta(K)$ permanently.

\paragraph{Relationship to the tractable-cognition literature.}
Work on probabilistic graphical models and treewidth analyzes the complexity of
inference \emph{within a fixed model structure} \cite{Pearl1988,vanRooij2008,Kwisthout2010}.
Our claims are different in kind: we focus on when representing shared
structure reduces the number of distinct quantities an agent must learn from $K$
to $F$, and on what revision transitions are required to recover that advantage
after model failure.

\begin{proposition}[Cost separation and recovery bound]\label{prop:crossover}
Let an agent face a $K$-option decision environment with $F$-factor structure
($F \leq K$).  Let $c_{\mathrm{param}}$ be the per-element parameter learning
cost and $c_{\mathrm{meta}}$ be the Level~2 model-revision overhead (navigation
of the revision graph $\mathcal{R}$, Section~\ref{sec:model_space}).  Then:

\begin{enumerate}[label=(\roman*)]
\item \textbf{Crossover.}
For $K < K^* = F + c_{\mathrm{meta}}/c_{\mathrm{param}}$, Linear achieves lower
total cognitive cost than Network (Network cannot do better than operating in a
correct compressed model and still paying its Level~2 overhead).
For $K > K^*$, Network under a correct factor hypothesis achieves lower total
cost by replacing $O(K)$ per-episode parameter maintenance with $O(F)$, up to the
fixed overhead $c_{\mathrm{meta}}$.

\item \textbf{Catastrophic failure under cross-terms.}
When the environment contains stakes-factor structure with covariance coupling,
Linear's utility degrades without bound in cumulative exposure $E_t$ because it
cannot represent the marginal penalty term $2E_t b_i$.  Network's utility
remains bounded because it tracks $E_t$ and computes marginal impact.

\item \textbf{Recovery bound.}
When the factor hypothesis is incorrect and the agent expands to
$G_{\mathrm{tab}}$, an agent whose revision graph $\mathcal{R}$ includes a
return transition $(G_{\mathrm{tab}}, G_F) \in \mathcal{T}$ recovers $O(F)$ cost
scaling within $n_{\mathrm{contract}} + w + n_{\mathrm{converge}}$ steps after
factor structure becomes present, where $n_{\mathrm{converge}}$ is the number of
observations required for within-factor variance to fall below the contraction
threshold.  An agent without a return transition (expand-only Network or Linear)
remains at $O(K)$ cost indefinitely.
\end{enumerate}

Parts (i) and (ii) are tested by the existing simulation results.
Part (iii) is tested by the contraction ablation
(Section~\ref{sec:contraction}).
\end{proposition}

\paragraph{Proof sketch for (i).}
In the absence of structural sharing, Linear must maintain $K$ parameters, so
its per-episode cost is proportional to $c_{\mathrm{param}}K$.  A Network agent
that is operating in a correct compressed model maintains only $F$ parameters,
but still pays a fixed Level~2 overhead $c_{\mathrm{meta}}$ for monitoring and
revision.  Comparing $c_{\mathrm{param}}K$ against $c_{\mathrm{param}}F +
c_{\mathrm{meta}}$ yields the crossover condition $K^* = F +
c_{\mathrm{meta}}/c_{\mathrm{param}}$.

\paragraph{Proof sketch for (ii).}
Under stakes, the true marginal impact of choosing option $i$ contains a
cross-term proportional to $2E_t b_i$, where $E_t$ is cumulative exposure.  A
Linear agent that treats options independently can include a per-option penalty
(e.g., $\lambda b_i^2$) but cannot represent the coupling to $E_t$, so it cannot
stabilize exposure as $|E_t|$ grows.  A Network agent that tracks $E_t$ and
computes marginal impact can keep exposure bounded, preventing unbounded penalty
growth.

\paragraph{Proof sketch for (iii).}
After expanding to $G_{\mathrm{tab}}$, an agent with a return transition can
periodically test whether a compressed factor model predicts well again (for
example via a length-$w$ evaluation window after a holdout period
$n_{\mathrm{contract}}$).  Once factor structure is present and sufficient data
are accumulated ($n_{\mathrm{converge}}$ observations), the agent contracts back
to $G_F$ and returns to $O(F)$ scaling.  Without a return transition, the agent
remains in the expanded model and continues to incur $O(K)$ cost.
